\documentclass[12pt,a4paper]{report}

\usepackage[utf8]{inputenc}
\usepackage[T1]{fontenc}
\usepackage[french]{babel}
\usepackage{geometry}
\usepackage{array}
\usepackage{booktabs}
\usepackage{longtable}
\usepackage{tabularx}
\usepackage{xcolor}
\usepackage{titlesec}
\usepackage{fancyhdr}
\usepackage{graphicx}
\usepackage{setspace}
\usepackage{enumitem}
\usepackage{hyperref}

\geometry{margin=2.5cm}
\setstretch{1.2}
\setlength{\parindent}{0pt}
\setlength{\parskip}{0.5em}

\definecolor{mainblue}{RGB}{26,87,132}
\definecolor{mainline}{RGB}{112,44,44}
\definecolor{softgray}{RGB}{245,245,245}
\definecolor{linegray}{RGB}{220,220,220}

\pagestyle{fancy}
\fancyhf{}
\fancyhead[C]{\textit{\textcolor{mainblue}{Fiche technique : Plateforme de laboratoire virtuel}}}
\fancyfoot[R]{\thepage}
\renewcommand{\headrulewidth}{0pt}

\titleformat{\chapter}[display]
  {\centering\bfseries}
  {}
  {0pt}
  {\Huge\itshape}

\titleformat{\section}
  {\large\bfseries}
  {\thesection.}
  {0.6em}
  {}

\titleformat{\subsection}
  {\normalsize\bfseries}
  {\thesubsection.}
  {0.6em}
  {}

\newcommand{\sectionline}{
  \vspace{0.3em}
  \noindent\color{mainline}\rule{\textwidth}{1.5pt}
  \vspace{0.6em}
}

\begin{document}

\begin{titlepage}
  \centering
  \vspace*{2.5cm}

  {\Huge\itshape Fiche technique\par}
  \vspace{0.8cm}
  {\Large\itshape de la plateforme de laboratoire virtuel\par}

  \vspace{2.2cm}

  {\Large\textbf{Projet :} EduVirtuel / Laboratoire virtuel interactif\par}
  \vspace{0.4cm}
  {\large Application web pédagogique avec intégration Unity WebGL\par}

  \vfill

  \begin{flushleft}
    \textbf{Contenu du document :}
    \begin{itemize}[leftmargin=1.2cm]
      \item Présentation générale du projet
      \item Architecture technique
      \item Outils et technologies utilisés
      \item Fonctionnalités réalisées
      \item Organisation des fichiers
      \item Travaux récents effectués
    \end{itemize}
  \end{flushleft}

  \vfill

  {\large Année universitaire : 2025--2026\par}
\end{titlepage}

\chapter{Plan de production et d'organisation}

\sectionline

\section{Présentation du projet}

Le projet réalisé est une plateforme de laboratoire virtuel destinée à un usage pédagogique.
Elle permet à l'élève d'accéder à plusieurs expériences interactives selon son niveau
scolaire, tout en offrant à l'enseignant un espace de gestion des classes, des codes
d'accès et du suivi des élèves.

La solution s'appuie sur deux composantes principales :

\begin{itemize}[leftmargin=1.2cm]
  \item une application web locale et déployable, construite en HTML, CSS et JavaScript ;
  \item un contenu immersif Unity WebGL réutilisé pour certains laboratoires 3D.
\end{itemize}

Le système comporte actuellement plusieurs laboratoires orientés vers la physique, les
circuits électriques, les sciences naturelles et la chimie, avec une interface élève et
une interface professeur.

\section{Objectif technique}

L'objectif de production a consisté à concevoir une plateforme simple à exécuter en local,
facile à maintenir, et adaptée à une démonstration pédagogique. Le projet devait permettre :

\begin{itemize}[leftmargin=1.2cm]
  \item l'accès aux laboratoires selon le niveau de l'élève ;
  \item la connexion et l'inscription des utilisateurs ;
  \item la gestion d'un tableau de bord enseignant ;
  \item l'intégration d'expériences interactives et d'un build Unity WebGL ;
  \item l'ajout d'un système de codes promotionnels de classe limitant certains laboratoires.
\end{itemize}

\section{Architecture générale de la solution}

\subsection{Partie web}

La partie principale du système est développée sous forme d'application web statique. Elle
comprend les pages d'accueil, les formulaires d'authentification, le tableau de bord élève,
le tableau de bord enseignant, les pages d'expériences et les pages de laboratoire.

\subsection{Partie simulation 3D}

Le projet réutilise également un build Unity WebGL ainsi qu'un projet Unity source plus ancien
destiné à la simulation de laboratoire virtuel. Cette partie complète l'expérience immersive
du système.

\subsection{Exécution locale}

Pour le lancement local, un serveur Python léger est utilisé afin de servir correctement les
fichiers WebGL, les ressources compressées \texttt{.gz} et les bons types MIME nécessaires.

\section{Organisation des fichiers du projet}

\sectionline

L'organisation du projet s'appuie principalement sur les dossiers suivants :

\begin{longtable}{|p{4cm}|p{10cm}|}
\hline
\rowcolor{softgray}
\textbf{Élément} & \textbf{Rôle} \\
\hline
\texttt{venerable-sprite-294b65-local/} & Application web principale du laboratoire virtuel. \\
\hline
\texttt{Laboratory-Simulator-sparse/} & Projet Unity source et build historique de simulation virtuelle. \\
\hline
\texttt{student.html} & Interface d'inscription et de connexion des élèves. \\
\hline
\texttt{prof-dashboard.html} & Tableau de bord enseignant pour le suivi et la génération des codes. \\
\hline
\texttt{script.js} & Logique principale de l'application : sessions, laboratoires, comptes, codes, tableaux de bord. \\
\hline
\texttt{style.css} & Mise en forme générale de l'interface. \\
\hline
\texttt{serve-local.py} & Serveur Python local adapté au chargement Unity WebGL. \\
\hline
\texttt{start-local.ps1} & Script PowerShell de démarrage rapide du projet en local. \\
\hline
\texttt{assets/unity-webgl/} & Build WebGL intégré dans la plateforme web. \\
\hline
\texttt{netlify.toml} & Configuration de déploiement de la partie web. \\
\hline
\end{longtable}

\section{Technologies et outils utilisés}

\sectionline

\subsection{Environnement logiciel}

Pour produire, modifier, tester et exécuter le projet, les outils suivants ont été utilisés
ou référencés dans le dépôt :

\begin{longtable}{|p{3.2cm}|p{3cm}|p{8.3cm}|}
\hline
\rowcolor{softgray}
\textbf{Outil / technologie} & \textbf{Version} & \textbf{Utilisation dans le projet} \\
\hline
Git & 2.43.0.windows.1 & Gestion de versions et suivi des modifications du projet. \\
\hline
Python & 3.11.4 & Exécution du serveur local \texttt{serve-local.py}. \\
\hline
Node.js & 22.14.0 & Vérifications rapides du code JavaScript et environnement de développement. \\
\hline
PowerShell & Environnement Windows & Lancement local, scripts et automatisation. \\
\hline
HTML5 & --- & Structure des pages de l'application web. \\
\hline
CSS3 & --- & Mise en forme graphique des interfaces. \\
\hline
JavaScript & --- & Logique métier, tableaux de bord, authentification locale et navigation. \\
\hline
Unity 3D & 2017.4.14f1 LTS & Moteur utilisé pour la simulation de laboratoire virtuelle historique. \\
\hline
Blender 3D & --- & Logiciel mentionné pour la modélisation 3D dans la partie Unity. \\
\hline
C\# & --- & Langage du projet Unity source. \\
\hline
Unity WebGL & --- & Export 3D intégré à la plateforme web. \\
\hline
Netlify & --- & Configuration de publication de la partie web. \\
\hline
\end{longtable}

\subsection{Bibliothèques et composants applicatifs}

La documentation du projet web mentionne également l'usage de bibliothèques spécialisées
pour certaines scènes interactives :

\begin{itemize}[leftmargin=1.2cm]
  \item \textbf{Three.js} pour le rendu 3D dans certaines vues web ;
  \item \textbf{cannon-es} pour la gestion de comportements physiques ;
  \item \textbf{GSAP} pour les animations et transitions de l'interface.
\end{itemize}

\section{Fonctionnalités réalisées}

\sectionline

Les principales fonctionnalités disponibles ou développées dans le cadre du projet sont les suivantes :

\begin{enumerate}[leftmargin=1.2cm]
  \item portail d'accès élève et enseignant ;
  \item inscription élève avec génération automatique d'un code personnel ;
  \item connexion ultérieure avec ce code personnel ;
  \item sélection et filtrage des expériences selon le niveau ;
  \item affichage d'expériences pour plusieurs domaines : physique, circuits, sciences, chimie ;
  \item tableau de bord enseignant avec suivi des comptes élèves ;
  \item création de codes d'accès par le professeur ;
  \item intégration de laboratoires 3D et d'un module Unity WebGL ;
  \item stockage local des données utilisateurs pour les démonstrations.
\end{enumerate}

\section{Travaux récents effectués sur le projet}

\sectionline

Durant la phase récente de développement, une amélioration importante a été apportée au flux
d'inscription et de gestion de classe.

\subsection{Ajout du système de code promo de classe}

Une nouvelle fonctionnalité a été développée afin de permettre au professeur de créer un
\textbf{code promo de classe}. Ce code ne correspond plus à une seule activité, mais à une
classe entière associée à un niveau donné.

Cette amélioration permet désormais :

\begin{itemize}[leftmargin=1.2cm]
  \item de créer un code de classe lié à un niveau (\textit{primaire}, \textit{CEM}, \textit{lycée}) ;
  \item d'autoriser uniquement certains laboratoires pour cette classe ;
  \item de faire entrer l'élève par un formulaire spécifique de type \textit{Code promo} ;
  \item de détecter automatiquement le niveau à partir du code fourni ;
  \item de créer ensuite un mot de passe personnel pour la connexion future de l'élève ;
  \item d'afficher uniquement les laboratoires autorisés par ce code dans l'espace élève.
\end{itemize}

\subsection{Compatibilité avec l'ancien fonctionnement}

Le nouveau système a été ajouté sans supprimer le comportement initial. Ainsi, deux modes
restent disponibles côté enseignant :

\begin{itemize}[leftmargin=1.2cm]
  \item \textbf{code d'activité unique} : ancien fonctionnement, limité à une seule activité ;
  \item \textbf{code promo de classe} : nouveau fonctionnement, avec plusieurs laboratoires autorisés.
\end{itemize}

\section{Organisation technique de la maintenance}

\sectionline

Pour faciliter la maintenance et l'évolution du projet, les règles suivantes ont été suivies :

\begin{itemize}[leftmargin=1.2cm]
  \item les pages HTML restent séparées selon les rôles et les vues ;
  \item le style global est centralisé dans \texttt{style.css} ;
  \item la logique applicative est regroupée dans \texttt{script.js} ;
  \item le serveur local est séparé dans \texttt{serve-local.py} ;
  \item les fichiers Unity WebGL sont isolés dans leur propre dossier d'assets ;
  \item la configuration de publication est séparée via \texttt{netlify.toml}.
\end{itemize}

\section{Conclusion}

Cette fiche technique montre que le projet repose sur une architecture hybride combinant
une interface web interactive, une logique JavaScript locale, un serveur Python de test,
et un contenu Unity WebGL pour les simulations immersives.

Le travail réalisé couvre non seulement la mise en place de l'application, mais aussi son
évolution fonctionnelle avec l'ajout d'un système de code promo de classe, mieux adapté à
un usage réel dans un contexte pédagogique.

\vspace{1em}

\textbf{Remarque :} ce document LaTeX est prêt à être compilé, mais l'outil \texttt{pdflatex}
n'est pas installé actuellement sur cette machine. Il est donc possible de générer le PDF
ultérieurement dès qu'une distribution LaTeX sera disponible.

\end{document}
